% ============================================================
% Abstract
% Included from main.tex
% Draft of 16 October 2026, written from the results measured so far;
% to be rewritten last, once the custom-set results exist.
% ============================================================

\chapter*{Abstract}
\addcontentsline{toc}{chapter}{Abstract}
% Mandatory (Art. 12). MDS requires it in English only (Art. 61);
% the trilingual abstract of Arts. 22, 28, 34 and 39 is for other
% degrees.

A publication of constructive news has to find, among many sources,
the few stories worth offering on the topics it covers, and check their
claims before it publishes them. This thesis builds and evaluates a
low-cost, end-to-end system that does both, in English and Spanish,
with open-weight models running locally and no paid service. It reads
37 configured news sources for up to three topic groups, screens out
what the publication never covers, ranks the candidates and lets an
editor, a language model or both choose at most twenty to analyse. For
each chosen article it extracts and enriches the text, admits or
rejects it, selects two to four checkable claims, searches the open web
for evidence through a self-hosted metasearch engine, ranks the
evidence with a gate that cuts off-target sources, and asks a local
language model for one of five verdicts that must cite and quote its
sources; the verdict's confidence is then recomputed from the evidence,
ungrounded verdicts are downgraded, and failed searches and unreachable
models are reported apart from verdicts.

The evaluation is the contribution. It is organised around four
research questions --- finding the news, verifying it, how close a
model comes to a person, and how much of a journalist's work the joined
system takes on --- and measures each stage against a person wherever a
person takes the same decision. It uses an evaluation harness that
reports every metric with a bootstrap interval and its provenance;
X-Fact as an external, open-retrieval benchmark; and a custom
validation set of 139 claims from constructive news (109 natural, 30
synthetic refuted claims built by substitution), the first of its
kind to the author's knowledge.

On a 64-claim X-Fact pilot the zero-shot 3-billion-parameter verifier
reached 29.7\% accuracy [18.8, 40.6] with a Cohen's $\kappa$ against
gold indistinguishable from zero, and 19.4\% [13.3, 25.5] on the 165
Spanish claims of X-Fact's in-domain test set; retrieved evidence added
10.9 points [3.1, 20.3] over the model alone on the pilot. The pilot
also showed how much of an open-retrieval score on an old benchmark
measures time rather than verification: most decidable claims cited
evidence newer than the claim, and the fact-checker's own verdict was
retrieved for over a third of the English ones. Before the verdict, the
ranked rounds separated stories worth offering from the rest with an
AUC of 0.78 [0.67, 0.87], and the claim selector's proposals were
checkable in 64\% of cases.
\pending{the custom-set results for RQ2--RQ4 with their intervals, the
agreement between the AI labelling pass and the annotator, and the
limitations that remain; then rewrite this abstract to one or two
pages.}

\textbf{Keywords:} automated fact-checking, open-domain retrieval,
large language models, constructive journalism, multilingual NLP,
evaluation.
